\documentclass[letterpaper,11pt]{article}
\usepackage[top=0.8in, bottom=0.8in, left=0.9in, right=0.9in]{geometry}
\usepackage[hidelinks]{hyperref}
\usepackage{lmodern}
\usepackage{parskip}
\usepackage{microtype}
\setlength{\parskip}{6pt}
\setlength{\parindent}{0pt}
\begin{document}
\begin{flushleft}
\textbf{\large Ajay Venkatesh} \\
Brooklyn, NY \\
+1 516 585 9013 \\
\href{mailto:av3855@nyu.edu}{av3855@nyu.edu}
\end{flushleft}
\vspace{10pt}
May 21, 2026
\vspace{6pt}

Hiring Manager \\
Appian Corporation

\vspace{10pt}

Dear Hiring Manager,

My name is Ajay Venkatesh. I'm finishing my M.S. in Computer Engineering at NYU, with production software experience at PwC in Bengaluru, a summer SDE internship at Amazon, and ongoing on-campus work at NYU's Novel AI Technologies. I'm writing about the Associate Application Engineer role.

The strongest match is the process-automation work I owned at PwC. I designed and delivered a \textbf{B2B onboarding platform} on \textbf{Microsoft Azure} integrating with \textbf{Microsoft Dynamics CRM} over \textbf{SQL Server}, modernizing a legacy operation by modeling business processes, building modular services, integrating data pipelines and email event tracking, and surfacing the results in BI reporting. The work was exactly the shape of process-centric application development described in this role: gather requirements with business stakeholders, model the workflow, build the integration layer, validate the deployment, and support production once it was live. I shipped that across multiple modules with strict review gates and a regulated PwC SDLC governance model.

In parallel, I built workflow-driven applications across the \textbf{Microsoft Power Platform} ecosystem (Power Automate, Power BI, Microsoft Copilot Studio, Dataverse), which is a close cousin of the \textbf{Appian platform} in spirit: low-code, business-process-first, integration-heavy. That hands-on experience translates directly into how I would think about implementing solutions on Appian, even though the platform itself would be net new to me.

At NYU I extended this into a customer-facing surface: I built a \textbf{broker portal} streamlining underwriting that integrated with external insurer platforms, with \textbf{document and media processing} and real-time communication. That gave me practice triaging production issues, supporting end users, and translating non-technical requirements into clean integrations.

On overlap with your stack: \textbf{AWS} (\textbf{Lambda, EC2, S3, DynamoDB, API Gateway, CDK}), \textbf{Azure} (Function Apps, AI Services, Logic Apps, Dataverse), \textbf{Microsoft Dynamics CRM}, \textbf{SQL Server / PostgreSQL / MySQL}, and \textbf{Python, Java, JavaScript / TypeScript}. I've also worked on \textbf{RAG} and \textbf{agentic} AI pipelines at Amazon and Campus Mesh, which lines up with the AI direction Appian is investing in. The \textbf{Appian platform itself} is the one piece I would need to learn, and I'd come up the curve quickly given how close it sits to the BPM and low-code work I've already done.

What pulls me toward Appian is the category. Process automation is one of the rare software domains where the work has direct, measurable consequences for how a business actually operates day to day, and a single well-modeled workflow can compound across hundreds of internal users. Combining that with the AI direction Appian is taking is exactly the kind of long-tail engineering work I would like to compound my career around.

I'd welcome the chance to discuss further. Thanks for considering my application.

\vspace{12pt}
Sincerely, \\
Ajay Venkatesh

\end{document}
